% Official CC BY4.0 published HTML S1; original math alttext and author text, literal original numbering.


\par\medskip\noindent\textbf{1 Introduction}\par




Strong computational complexity lower bounds are often difficult to prove, for many different complexity measures. In most cases, the inability to prove lower bounds can be explained by barrier results such as the Relativization barrier due to Baker, Gill and Solovay [5], the Algebrization barrier due to Aaronson and Wigderson [1] and the Natural Proofs barrier due to Razborov and Rudich [66]. Metacomplexity studies the complexity of computing various complexity measures, such as the minimum circuit complexity of a function.
One eventual goal of the metacomplexity program is to understand mathematically why lower bounds are hard to prove.
In this paper, we study metacomplexity in the algebraic circuit complexity setup from a representation-theoretic viewpoint.





Algebraic complexity focuses on proving complexity lower bounds for various problems related to the computation of multivariate polynomials and rational functions. Prominent examples include the discrete Fourier transform, matrix multiplication, and solving systems of linear equations; see [14].
In algebraic complexity, the main model of computation is the algebraic circuit, which is a labeled directed graph encoding an algorithm to evaluate a polynomial. Most polynomials require large algebraic circuits; however, exhibiting explicit ones with such property is difficult. Valiant’s flagship conjecture [71], known as $\mathsf{VP}\neq\mathsf{VNP}$, states that the algebraic circuit complexity of the permanent polynomial $\mathsf{per}_{d}(x_{1,1},\dots,x_{d,d})=\sum_{\pi\in\mathfrak{S}_{d}}\prod_{i=1}^{d}x_{i,\pi(i)}$ is not polynomially bounded as a function of $d$. This conjecture is an algebraic analogue of the $\mathsf{P}\neq\mathsf{NP}$ conjecture: indeed, the permanent of the adjacency matrix of a bipartite graph counts the number of perfect matchings, which is a $\#\mathsf{P}$-hard problem.





Many current proof techniques in algebraic complexity theory are implicitly based on metapolynomials. Examples of such techniques are rank methods, dating back to [70], and used more recently in complexity theory both for polynomials [62, 55] and for tensors [69, 54], methods based on geometric features [53, 36], and the geometric complexity program [60, 59]. We refer to Section 2.2 for a more extensive discussion. Informally, metapolynomials are polynomials whose variables are coefficients of polynomials. The concept dates back to resultants and invariants of forms [20, 70], while the term “metapolynomial” was coined in [67] and made popular in [40]. A classical example is the metapolynomial $b^{2}-4ac$ in the coefficients of a univariate degree $2$ polynomial $ax^{2}+bx+c$, which vanishes exactly when the two roots of the polynomial coincide.
Grochow pointed out in [39] that for many border complexity measures the complexity lower bounds can without loss of generality be proved by metapolynomials that lie in separating modules, or are isotypic, or are highest weight vectors.
These are all closely related concepts (and we prove analogous results for all of them), where being isotypic is the most natural one because it does not depend on choosing any bases.
Being a homogeneous metapolynomial is the same as being isotypic with respect to the action of $1\times 1$ matrices.
Being isotypic also generalizes the property of being invariant under the special linear group.
While invariants can have large circuit size, see [34],
we prove that not all separating isotypic metapolynomials have large circuit size, unless all separating metapolynomials have large circuit size.
More precisely, we prove that without loss of generality one may assume the metapolynomial to be isotypic, or to be a highest weight vector, or to lie in a separating module, with only a quasipolynomial blowup in the circuit size.
We use this to answer Open Question 2 in [40] by proving that all metapolynomials that lie in an irreducible representation have the same circuit complexity up to quasipolynomial blowup, see Corollary 1.2.





Our result reduces the search space for lower bounds from all homogeneous metapolynomials to the isotypic metapolynomials, which is a technique for example used in
[2, 15, 16, 43, 3, 23, 11]. In particular the large-scale cluster computations in [3, 23, 11] would be infeasible without this search space reduction. For example, [43] reduces the search space for proving that the border complexity of the $2\times 2$ matrix multiplication tensor is 7 to a 4-dimensional linear subspace of isotypic metapolynomials in a $\binom{64+20-1}{20}=8,179,808,679,272,664,720$-dimensional space of homogeneous degree 20 metapolynomials on $4\times 4\times 4$ tensors. On the reduced search space, the problem is then easily solved by linear algebra.
This specific approach was used in [8] to prove that the border support rank of the $2\times 2$ matrix multiplication tensor is 7. To this date, this is the only proof of this fact.
Similar methods are used in [63, 37], where the polynomials of interest are invariants for the action of a special linear group: this information is crucially used to reduce the search space and determine the polynomial via interpolation methods on a small space.
These are finite results, but the techniques work in the asymptotic settings as well, for example for lower bounds on the rank of matrix multiplication as in [15, 16], or the
best known lower bound for the border determinantal complexity of the permanent
[53].





Our result also points to potential future lower bounds:
Our method can be used to extract isotypic metapolynomials for example from rank-based lower bounds proofs, and the obtained metapolynomials give at least the same lower bounds, but potentially better ones.
It is conceivable that these metapolynomials can be adjusted to not be affected by barriers for rank methods [26, 33].





Our procedure uses explicit efficient circuit constructions and hence makes a contribution to algorithmic representation theory that is of independent interest.





\par\medskip\noindent\textbf{Metapolynomials}\par




We work over the field $\mathbb{C}$ of complex numbers, and the set $\{x_{1},x_{2},\ldots\}$ of variables.
For every sequence $\bm{i}$ of nonnegative integers with $\sum_{j}i_{j}<\infty$,
a monomial $x^{\bm{i}}$ is the finite product of variables $x_{1}^{i_{1}}x_{2}^{i_{2}}\cdots$, and a polynomial $f$ is a finite linear combination of monomials.
For every monomial $x^{\bm{i}}=x_{1}^{i_{1}}x_{2}^{i_{2}}\cdots$, define a metavariable $c_{\bm{i}}$.
A metavariable $c_{{\bm{i}}}$ has degree 1, and it has weight $\bm{i}$.
Finite products of metavariables are called metamonomials, and finite linear combinations of metamonomials are called metapolynomials.
The weight of a metamonomial is the sum of the weights of its metavariables.
If every metamonomial in a metapolynomial $\Delta$ has the same weight, then $\Delta$ is called a weight metapolynomial (and the zero metapolynomial has every weight).
For a fixed weight $w$ the metapolynomials of that weight form a vector space, the weight $w$ space.
Every metapolynomial can be written as a unique sum of weight metapolynomials of different weights.
A metapolynomial is called affine linear if all its metamonomials have degree at most 1, i.e., are just a metavariable or a constant.
Every metapolynomial $\Delta$ can be evaluated at any polynomial $f$ by assigning to each metavariable the coefficient of the corresponding monomial in $f$.
For example, the discriminant of a degree two homogeneous polynomial in two variables $c_{20}x_{0}^{2}+c_{11}x_{0}x_{1}+c_{02}x_{1}^{2}$ is the metapolynomial $\Delta=c_{11}^{2}-4c_{20}c_{02}$. It is a classical fact that $\Delta(f)=0$ if and only if $f=\ell^{2}$ for some homogeneous linear polynomial $\ell=\alpha_{0}x_{0}+\alpha_{1}x_{1}$.
In principle, metapolynomials can involve metavariables that correspond to monomials of different degrees, but for our purposes it is sufficient to consider the case where all metavariables correspond to monomials of the same degree $d$, see Remark 1.3.
Every such metapolynomial can be decomposed uniquely as a sum of its homogeneous degree $\delta$ components,
and we call a homogeneous degree $\delta$ metapolynomial a metapolynomial of format $(\delta,d,k)$ if all its metavariables only involve the variables $x_{1},\ldots,x_{k}$.






\par\medskip\noindent\textbf{Algebraic circuits}\par




In this paper, we are mostly interested in the algebraic circuit size of metapolynomials. Algebraic circuits for metapolynomials are defined in the same way as for usual polynomials, as follows.





An algebraic circuit is a directed acyclic graph in which each vertex
has indegree $0$ or $2$. Every indegree 0 vertex is called input gate and it is labelled by an affine linear metapolynomial; every indegree 2 vertex is called a computation gate and it is labeled by either “$+$” or “$\times$”; every edge is labeled by a complex number, which is assumed to be 1 if this scalar is omitted; there is exactly one vertex of outdegree 0. An algebraic circuit computes a metapolynomial at every vertex by induction over the directed acyclic graph structure (the labels on the edges rescale the computed values), i.e.,
if the values computed at the children of a gate $v$ are $\Delta_{1}$ and $\Delta_{2}$ and the incoming edges to $v$ are labelled with $\alpha_{1}$ and $\alpha_{2}$, respectively, then the value computed at $v$ is $\alpha_{1}\Delta_{1}+\alpha_{2}\Delta_{2}$ or $\alpha_{1}\alpha_{2}\Delta_{1}\Delta_{2}$, depending on whether $v$ is labelled with a “$+$” or a “$\times$”.
An algebraic circuit computes the metapolynomial computed at its outdegree 0 vertex.
The size of a circuit is the total number of its vertices. The algebraic circuit complexity of a metapolynomial is the minimum size of an algebraic circuit computing it.
For example, Figure 1 shows that the algebraic circuit complexity of the discriminant is at most 6. Algebraic circuits for polynomials are defined analogously.




Figure 1: An algebraic circuit computing the discriminant $\Delta=c_{11}^{2}-4c_{20}c_{02}$.




\par\medskip\noindent\textbf{Representation Theory}\par




Let $\mathbb{N}$ denote the set of natural numbers including zero.
A partition $\lambda=(\lambda_{1},\lambda_{2},\ldots)$ is a finite sequence of non-increasing natural numbers.
Let $\ell(\lambda)=\max\{j\mid\lambda_{j}\neq 0\}$. We define $\lambda_{j}=0$ for all $j>\ell(\lambda)$.
Write $|\lambda|:=\sum_{j}\lambda_{j}$, and say that $\lambda$ is a partition of $|\lambda|$ of length $\ell(\lambda)$. In this case we write $\lambda\mathchoice{\scalebox{1.2}[0.7]{ $\vdash$}\,}{\scalebox{1.2}[0.7]{ $\vdash$}\,}{\scalebox{0.7}{\scalebox{1.2}[0.7]{$\vdash$}}\,}{\scalebox{0.5}{\scalebox{1.2}[0.7]{$\vdash$}}\,}|\lambda|$ or $\lambda\mathchoice{\scalebox{1.2}[0.7]{ $\vdash$}_{\ell(\lambda)}\,}{\scalebox{1.2}[0.7]{ $\vdash$}_{\ell(\lambda)}\,}{\scalebox{0.7}{\scalebox{1.2}[0.7]{$\vdash$}}_{\ell(\lambda)}\,}{\scalebox{0.5}{\scalebox{1.2}[0.7]{$\vdash$}}_{\ell(\lambda)}\,}|\lambda|$ depending on whether the length is relevant. For example, $\lambda=(4,2,2)$ is a partition of $8$ of length $3$, i.e., $\lambda\mathchoice{\scalebox{1.2}[0.7]{ $\vdash$}_{3}\,}{\scalebox{1.2}[0.7]{ $\vdash$}_{3}\,}{\scalebox{0.7}{\scalebox{1.2}[0.7]{$\vdash$}}_{3}\,}{\scalebox{0.5}{\scalebox{1.2}[0.7]{$\vdash$}}_{3}\,}8$.
The Young diagram of a partition $\lambda$ is a top-left justified array of boxes with $\lambda_{j}$ many boxes in row $j$. For example, the Young diagram for $(4,2,2)$ is 
. A Young tableau $T$ is a filling of the boxes of a Young diagram with positive integer numbers. The partition $\lambda$ is called the shape of $T$. For example, 
is a Young tableau of shape $(4,2,2)$. A Young tableau is semistandard if the numbers within each row from left to right are non-decreasing, and the numbers within each column from top to bottom are strictly increasing, for example 
.
For each partition $\lambda$ there exists a unique superstandard tableau, that is, the tableau with only the number $j$ in row $j$, e.g., 
.





The group $\textnormal{GL}_{k}$ acts on the space $\mathbb{C}[x_{1},\ldots,x_{k}]_{d}$ of homogeneous polynomials on $\mathbb{C}^{k}$ of degree $d$ by linear change of coordinates. This action induces, again by linear change of coordinates, an action on the space of metapolynomials with format $(\delta,d,k)$ for every $\delta$.
For instance, the change of basis swapping the two coordinates on $\mathbb{C}^{2}$, defines the change of basis on $\mathbb{C}[x_{0},x_{1}]_{2}$ obtained by exchanging $x_{0}$ and $x_{1}$. This induces a linear change of coordinates on $\mathbb{C}[c_{20},c_{11},c_{02}]$ which exchanges $c_{20}$ and $c_{02}$ and leaves $c_{11}$ fixed. The discriminant $c_{11}^{2}-4c_{02}c_{20}$ is sent to itself by this change of coordinates.
More precisely the group actions are defined as follows. For every homogeneous degree $d$ polynomial $f$ in $k$ variables and every $g\in\textnormal{GL}_{k}$, the polynomial $g\cdot f$ is defined by $(g\cdot f)(x)=f(g^{-1}x)$ for every $x\in\mathbb{C}^{k}$. Similarly, for every format $(\delta,d,k)$ metapolynomial $\Delta$, the metapolynomial $g\cdot\Delta$ is defined by $(g\cdot\Delta)(f)=\Delta(g^{-1}\cdot f)$ for every homogeneous degree $d$ polynomial $f$. The details about this action are provided in Section 4.





The vector space of metapolynomials of format $(\delta,d,k)$ is closed under this action of $\textnormal{GL}_{k}$, and it decomposes into a direct sum of subspaces which are also closed under the action of $\textnormal{GL}_{k}$: these subspaces are called isotypic components and they are in one to one correspondence with partitions $\lambda\mathchoice{\scalebox{1.2}[0.7]{ $\vdash$}_{k}\,}{\scalebox{1.2}[0.7]{ $\vdash$}_{k}\,}{\scalebox{0.7}{\scalebox{1.2}[0.7]{$\vdash$}}_{k}\,}{\scalebox{0.5}{\scalebox{1.2}[0.7]{$\vdash$}}_{k}\,}{\delta d}$.
The summand corresponding to the partition $\lambda$ is called the $\lambda$-isotypic component. This is called the isotypic decomposition. Each $\lambda$-isotypic component decomposes even further into a direct sum of subspaces with one summand for each semistandard tableau $T$ of shape $\lambda$.
We call the component for $T$ the $T$-isotypic component.
If $V$ is $\lambda$-isotypic and irreducible (i.e., has no nontrivial subrepresentations), then $\lambda$ is called the isomorphism type of $V$, and each $T$-isotypic component of $V$ is 1-dimensional.
If $T$ is a superstandard tableau of shape $\lambda$, then the corresponding $T$-isotypic component is called the highest weight metapolynomial vector space of weight $\lambda$.
Equivalently, a highest weight vector of weight $\lambda$ is a metapolynomial $\Delta$ with $g\cdot\Delta=t_{1}^{\lambda_{1}}\cdots t_{k}^{\lambda_{k}}\Delta$, for every $g\in\textnormal{GL}_{k}$ upper triangular with $t_{1},\ldots,t_{k}$ on the main diagonal. The $\lambda$-isotypic component is known to be the linear span of the union of the orbits of all its highest weight vectors.





For example (not easily verifiable by hand, see Example 5.4), the space of metapolynomials of format $(\delta,d,k)=(3,2,3)$ decomposes into three nonzero isotypic components, corresponding to partitions $(6,0,0)$, $(4,2,0)$ and $(2,2,2)$ of $\delta\cdot d=6$. The isotypic component $(4,2,0)$ further decomposes into three $T$-isotypic components for the three semistandard tableaux of shape $(4,2,0)$: 
, 
, 
. The metapolynomial $\Delta=c_{200}c_{020}c_{002}$ can be written according to this decomposition as



\[\begin{aligned}
\displaystyle 60\,\Delta & \displaystyle= \displaystyle 2\big{(}c_{101}^{2}c_{020}+2c_{110}c_{101}c_{011}+c_{200}c_{011}^{2}+c_{110}^{2}c_{002}+2c_{200}c_{020}c_{002}\big{)} \\
 & \displaystyle+ \displaystyle\Big{[}\ 2\big{(}-c_{020}c_{101}^{2}-2c_{011}c_{101}c_{110}+4c_{002}c_{110}^{2}-c_{011}^{2}c_{200}+8c_{002}c_{020}c_{200}\big{)} \\
 & \displaystyle+ \displaystyle\phantom{\Big{[}\ }0 \\
 & \displaystyle+ \displaystyle\phantom{\Big{[}\ }5\big{(}c_{020}c_{101}^{2}-c_{011}c_{101}c_{110}-c_{002}c_{110}^{2}+c_{011}^{2}c_{200}+4c_{002}c_{020}c_{200}\big{)}\ \Big{]} \\
 & \displaystyle+ \displaystyle 5\big{(}-c_{101}^{2}c_{020}+c_{110}c_{101}c_{011}-c_{200}c_{011}^{2}-c_{110}^{2}c_{002}+4c_{200}c_{020}c_{002}\big{)},
\end{aligned}\]


where the five summands from top to bottom are 
-isotypic,

-isotypic,

-isotypic,

-isotypic,
and

-isotypic, respectively.
The fifth summand is a highest weight vector, because 
is superstandard.
All summands have weight (2,2,2), because $\Delta$ has weight (2,2,2).





\par\medskip\noindent\textbf{Theorem 1.1 (Main theorem).}\par



Let $\Delta\colon\mathbb{C}[x_{1},\dots,x_{k}]_{d}\to\mathbb{C}$ be a metapolynomial of format $(\delta,d,k)$ computed by an algebraic circuit of size $s$.
Then



\begin{enumerate}


\item 
 


For every weight $\mu\in\mathbb{N}^{k}$, $|\mu|=d\delta$, the projection of $\Delta$ onto the weight space of weight $\mu$ can be computed by a circuit of size $O(s(\delta d)^{2k^{3}})$.





\item 
 


For every partition $\lambda\mathchoice{\scalebox{1.2}[0.7]{ $\vdash$}\,}{\scalebox{1.2}[0.7]{ $\vdash$}\,}{\scalebox{0.7}{\scalebox{1.2}[0.7]{$\vdash$}}\,}{\scalebox{0.5}{\scalebox{1.2}[0.7]{$\vdash$}}\,}d\delta$ the projection of $\Delta$ onto the $\lambda$-isotypic component can be computed by a circuit of size $O(sk^{2k^{2}}(\delta d)^{2k^{3}})$.





\item 
 


For every partition $\lambda\mathchoice{\scalebox{1.2}[0.7]{ $\vdash$}\,}{\scalebox{1.2}[0.7]{ $\vdash$}\,}{\scalebox{0.7}{\scalebox{1.2}[0.7]{$\vdash$}}\,}{\scalebox{0.5}{\scalebox{1.2}[0.7]{$\vdash$}}\,}d\delta$ the projection of $\Delta$ onto the highest weight space of weight $\lambda$ can be computed by a circuit of size $O(s(k+1)^{2k^{2}}(\delta d)^{2k^{3}})$.





\item 
 


For every semistandard tableau $T$ of shape $\lambda\mathchoice{\scalebox{1.2}[0.7]{ $\vdash$}\,}{\scalebox{1.2}[0.7]{ $\vdash$}\,}{\scalebox{0.7}{\scalebox{1.2}[0.7]{$\vdash$}}\,}{\scalebox{0.5}{\scalebox{1.2}[0.7]{$\vdash$}}\,}d\delta$ the projection of $\Delta$ onto the $T$-isotypic space can be computed by a circuit of size $O(sk^{2k^{2}}(\delta d)^{2k^{4}})$.





\end{enumerate}





In all applications of
Theorem 1.1,
we will always have that $k$ is logarithmic in the number of metavariables, which we explain in Section 2.5.
Open Question 2 in [40] asks: given a sequence of irreducible representations of metapolynomials, what is the sequence of metapolynomials in it that has the lowest circuit complexity? Theorem 1.1 answers this question in a very satisfying way: Every such sequence of metapolynomials has the same circuit complexity up to quasipolynomial blowup, as seen in the following corollary.





\par\medskip\noindent\textbf{Corollary 1.2.}\par



Let $s$ be the smallest circuit complexity of a nonzero metapolynomial in an irreducible $\textnormal{GL}_{k}$-representation $V$ of type $\lambda\mathchoice{\scalebox{1.2}[0.7]{ $\vdash$}\,}{\scalebox{1.2}[0.7]{ $\vdash$}\,}{\scalebox{0.7}{\scalebox{1.2}[0.7]{$\vdash$}}\,}{\scalebox{0.5}{\scalebox{1.2}[0.7]{$\vdash$}}\,}\delta d$.
Then every metapolynomial in $V$ has circuit complexity at most
$O(sk^{2k^{2}}(\delta d)^{2k^{4}+k^{2}})$.






\par\medskip\noindent\textbf{Proof.}\par




If $V$ is irreducible of isomorphism type $\lambda$, then it decomposes into a direct sum of $1$-dimensional $T$-isotypic components, where $T$ ranges over all semistandard tableaux of shape $\lambda$.
Take a nonzero metapolynomial $\Delta$ with minimal complexity in $V$ and apply a generic basis change, i.e., apply a generic element of $\textnormal{GL}_{k}$.
Then independently project the result onto all $T$-isotypic components.
We claim that since the basis change was generic, the projections are nonzero, and therefore form a basis of $V$.
In fact, we can be more precise about how we pick the basis change element $g\in\textnormal{GL}_{k}$.
For every semistandard tableau $T$, define $X_{T}:=\{g\in\textnormal{GL}_{k}\mid g\Delta\text{ has zero $T$-isotypic projection}\}$.
Such $X_{T}$ is Zariski closed in $\textnormal{GL}_{k}$, but it is not the whole $\textnormal{GL}_{k}$, otherwise the linear span of $\textnormal{GL}_{k}\Delta$ would be a nontrivial subrepresentation of the irreducible representation $V$.
Therefore the finite union $X:=\bigcup_{T}X_{T}$ is also a proper Zariski closed subset of $V$, and its complement $\textnormal{GL}_{k}\setminus X$ is a nonempty Zariski open subset.
If we pick $g\in\textnormal{GL}_{k}\setminus X$, then the projections of $g\Delta$ to every $T$-isotypic space are nonzero.
Moreover, by Theorem 1.1, all these projections have complexity at most $O(sk^{2k^{2}}(\delta d)^{2k^{4}})$.
But the dimension of $V$ equals the number of semistandard tableaux of shape $\lambda$ with entries $1,\ldots,k$, which is at most $(\delta d)^{k^{2}}$ (each tableaux has $k$ rows, and each row has length at most $\delta d$, so can be filled in at most $(\delta d)^{k}$ ways).
$\hfill\blacktriangleleft$





\par\medskip\noindent\textbf{$\blacktriangleright$ Remark 1.3.}\par



In general, algebraic complexity classes are defined for possibly nonhomogeneous polynomials. In this paper, we assume that polynomials and metapolynomials are homogeneous. However, one can lift the results to the nonhomogeneous setting as follows.





A polynomial $f(x_{1},\ldots,x_{k})$ is called homogeneous if all its monomials have the same degree.
For $d\geq\deg(f)$, define the degree $d$ homogenization



\[\begin{aligned}
f^{\sharp d}(x_{0},x_{1},\ldots,x_{k}):=x_{0}^{d}f(x_{1}/x_{0},\ldots,x_{k}/x_{0}).
\end{aligned}\]


If $f$ has an algebraic circuit of size $s$, then $f^{\sharp d}$ has an algebraic circuit of size at most $O(sd)$:
$s(d+1)$ for extracting the homogeneous parts of $f$ via interpolation,
$d$ operations to compute all values $x_{0},x_{0}^{2},\ldots,x_{0}^{d}$,
and another $d+1$ for multiplying the homogeneous components with the correct powers of $x_{0}$, and a final $d$ for adding up the results.
Given a metapolynomial $\Delta$ that does not involve $x_{0}$,
define $\Delta^{\sharp d}$ by replacing every metavariable $c_{\bm{i}}$ with $c_{(d-|\bm{i}|,i_{1},i_{2},\ldots)}$,
and setting every metavariable $c_{\bm{i}}$ with $|\bm{i}|>d$ to zero.
Clearly we have $\Delta^{\sharp d}(f^{\sharp d})=\Delta(f)$, provided that $d\geq\deg(f)$.
Hence, in this paper we can always assume that $f$ is homogeneous.








\par\medskip\noindent\textbf{Organization of the paper}\par




In the first part of Section 2 we illustrate the motivations for studying the circuit complexity of the isotypic components and the highest weight vectors in the space of metapolynomials, in the context of algebraic complexity lower bounds and border complexity classes.
In Section 2.5, based on Theorem 1.1 we prove that algebraic natural proofs can without loss of generality be assumed to be isotypic, see Theorem 2.4.
Section 3 sketches the main ideas behind Theorem 1.1. In Section 4, we provide an introduction to Lie algebras, universal enveloping algebras, and their representation theory, on which the core of the proof is built. In Section 5.1, we explain the construction of the projection maps used in the proof of the main theorem; we describe the efficient circuit implementation in Section 5.2. In Appendix A, we prove the results of Section 2.5. In Appendix B, we show one part of the proof of the Poincaré–Birkhoff–Witt Theorem (Theorem 4.9); this proof is explicit and can be used to construct the projectors in Theorem 1.1.




